\section{谱性公理}
\label{sec:spectrality-axiom}

一般辛流形、特别是无理辛流形 $(M,\omega)$ 上的谱不变量 $\rho(H;a)$ 有如下最为
非平凡的性质之一：
\begin{equation}
 \text{（谱性）}\qquad \rho(H;a)\in\operatorname{Spec}(H)
 \label{eq:21-spectrality}
\end{equation}
对\emph{任意} $H$ 与 $a\in QH^*(M)$ 成立。对任意\emph{有理}辛流形，即使 $H$ 退化，这一性质
也相对容易证明；见 (Oh05c)。
\footnote{译者注：原文此处再次把谱性写成对所有 $a\in QH^*(M)$ 成立，但
定义~\ref{def:21-spectral-invariant} 只对 $a\neq0$ 定义实值谱不变量；译文照录。}

另一方面，对\emph{无理}辛流形迄今只证明了以下较弱版本；退化哈密顿量是否满足谱性仍是
未知的：
\stfhIndex{I136}{Feituihuapuxing}{非退化谱性（nondegenerate spectrality）}
\[
 \tag{非退化谱性}
 \text{对任意非退化 }H\text{ 及 }a\in QH^*(M),
 \quad\eqref{eq:21-spectrality}\text{ 成立。}
\]

\begin{question}
\label{prob:21-degenerate-irrational-spectrality}
无理辛流形上的退化哈密顿量是否满足谱性公理？若否，构造一个违反谱性公理的
哈密顿量。
\end{question}

不过，结合柔性逼近论证，这一非退化谱性足以证明一般的 $\rho(H;a)$ 同伦不变性；
\stfhIndex{I169}{JinFuluoerbilian}{紧弗洛尔闭链（tight Floer cycle）}
见 21.6 节。(Oh05c, Oh09a) 给出的谱性证明方案依赖\emph{紧弗洛尔闭链}的存在性。

\begin{definition}
\label{def:21-tight-floer-cycle}
若弗洛尔闭链 $\alpha\in CF(H)$ 满足：对任意与其同调的弗洛尔闭链
$\alpha'\in CF(H)$，即 $\alpha'-\alpha=\partial_{(H,J)}\gamma$，都有
\[
 \lambda_H(\alpha')\geq\lambda_H(\alpha),
\]
则称 $\alpha$ 是\emph{紧}的。
\end{definition}

本节尝试在起初不假设辛流形有理性的前提下证明谱性，解释这个几何论证在无理情形
下可能出错之处，再对任意哈密顿量——无论是否非退化——证明有理情形的谱性。
关于无理情形非退化谱性的完整证明，参见 (Oh09a) 或 (Ush08)。(Oh09a) 的证明是
几何的，并在称作\emph{瑟夫同伦}的通有单参数哈密顿量族下证明谱不变量的许多其他性质；
(Ush08) 的证明则是纯代数的，因为它适用于 13.8 节意义下的任意抽象弗洛尔复形。

\stfhIndex{I050}{Dianfanbozhumian}{典范薄柱面（canonical thin cylinder）}
回忆\termCanonicalThinCylinder{} (18.2.7) 的定义。重要的是，$u_{zz'}^{\mathrm{can}}$ 的像包含在
$z$（或 $z'$）的小邻域内；当 $z,z'$ 在 $C^1$ 拓扑中趋于环路 $z_\infty$ 时，它一致
收敛到 $z_\infty$。因此，在满足边界条件
\[
 u(0,t)=z(t),\qquad u(1,t)=z'(t)
\]
的映射 $u:[0,1]\times S^1\to M$ 的同伦类集合中，$u_{zz'}^{\mathrm{can}}$ 也选出
一个典范同伦类 $[u_{zz'}^{\mathrm{can}}]$。

\begin{lemma}
\label{lem:21-canonical-cylinder-area}
设 $z,z':S^1\to M$ 是两条光滑环路，$u^{\mathrm{can}}$ 是 (18.2.7) 给出的上述典范
圆柱。对 $z$ 的任意填充圆盘 $w$ 及附近环路 $z'$，关联 $z'$ 的填充圆盘
\[
 w':=w\#u_{zz'}^{\mathrm{can}}.
\]
当 $d_{C^1}(z,z')\to0$ 时，$u_{zz'}^{\mathrm{can}}$ 在 $C^1$ 拓扑中收敛，且其几何
面积 $\operatorname{Area}(u^{\mathrm{can}})$ 趋于零。特别地，当 $z'\to z$ 且
$d_{C^1}(z,z')\to0$ 时：
\footnote{译者注：原文只说 $u_{zz'}^{\mathrm{can}}$ 在 $C^1$ 拓扑中“收敛”，
却没有给出其极限；随后第 (1) 项明确陈述的是拼接圆盘 $w'\to w$。译文照录。}
\begin{enumerate}[label=(\arabic*)]
 \item 作为单位圆盘到 $M$ 的映射，$w'\to w$ 于 $C^1$ 拓扑；
 \item 辛面积满足
 \begin{equation}
  \int_{u_{zz'}^{\mathrm{can}}}\omega\longrightarrow0.
  \label{eq:21-thin-cylinder-symplectic-area}
 \end{equation}
\end{enumerate}
\end{lemma}

\begin{proof}
第 (1) 项立即来自上述\termCanonicalThinCylinder{}的显式形式及指数映射的标准性质。另一方面，由\termCanonicalThinCylinder{}的显式表达与指数映射的性质，容易作面积估计，得到当
$d_{C^1}(z,z')\to0$ 时 $\operatorname{Area}(u_{zz'}^{\mathrm{can}})\to0$。因
$z,z'$ 均为 $C^1$ 的，$u_{zz'}^{\mathrm{can}}$ 也是 $C^1$ 的，故
\[
 \operatorname{Area}(u_{zz'}^{\mathrm{can}})
 \geq\left|\int_{u_{zz'}^{\mathrm{can}}}\omega\right|.
\]
这蕴含当 $z'\to z$ 于 $C^1$ 拓扑时，
$\int_{u_{zz'}^{\mathrm{can}}}\omega\to0$。
\end{proof}

目标是证明极小极大值 $\rho(H;a)$ 是临界值，即存在
$[z,w]\in\widetilde{\mathcal L}_0(M)$ 使
\[
 \mathcal A_H([z,w])=\rho(H;a),\qquad
 d\mathcal A_H([z,w])=0,
 \quad\text{即}\quad \dot z=X_H(z).
\]
此前已证明 $\rho(H;a)$ 有限。

\begin{theorem}
\label{thm:21-rational-spectrality}
设 $(M,\omega)$ 有理，则任意哈密顿量 $H$ 均满足谱性公理，不要求 $H$ 非退化。
\end{theorem}

\begin{proof}
我们可以通过反证法证明：对于任意光滑哈密顿量，谱性公理成立，证明如下。
若 $H$ 非退化，就固定 $H$；否则在 $C^2$ 拓扑中以非退化哈密顿量序列
$H_i$ 逼近 $H$。选取弗洛尔闭链 $\alpha_i\in CF_*(H_i)$，使
$\lambda_{H_i}(\alpha_i)\to\rho(H;a)$。令
$\operatorname{peak}(\alpha_i)=[z_i,w_i]\in\operatorname{Crit}\mathcal A_{H_i}$。
则
\begin{equation}
 \lim_{i\to\infty}\mathcal A_{H_i}([z_i,w_i])=\rho(H;a).
 \label{eq:21-peak-action-limit}
\end{equation}
由 $\rho(\,\cdot\,;a)$ 的定义及有限性，可选出这样的序列。

现在我们证明序列 $[z_i,w_i]$ 在 $\widetilde{\mathcal L}_0(M)$ 上是预紧的。
因 $M$ 紧、$H_i\to H$ 于 $C^2$ 拓扑且 $\dot z_i=X_{H_i}(z_i)$，标准自举论证给出
$z_i$ 的一个子列（仍记 $z_i$）收敛到某环路 $z_\infty:S^1\to M$，且
$\dot z_\infty=X_H(z_\infty)$。由于固定了非零量子上同调类
$a\in QH^*(M)$（更具体地说，固定了它的次数），并且 $[\alpha_i]=a^\flat$，有
\[
 \mu_{H_i}([z_i,w_i])=\mu_{H_j}([z_j,w_j]).
\]

对充分大的 $i$，考虑 $z_\infty$ 的填充圆盘
\[
 w'_i=w_i\#u_{i\infty}^{\mathrm{can}},
\]
其中 $u_{i\infty}^{\mathrm{can}}=u_{z_i z_\infty}^{\mathrm{can}}$ 是 $z_i$ 与
$z_\infty$ 之间的\termCanonicalThinCylinder{}。其几何面积在 $i\to\infty$ 时趋于零。计算临界点
$[z_\infty,w'_i]\in\operatorname{Crit}\mathcal A_H$ 的作用量：
\begin{align}
 \mathcal A_H([z_\infty,w'_i])
 &=-\int_{w'_i}\omega-\int_0^1H(t,z_\infty(t))\,dt \notag\\
 &=-\int_{w_i}\omega-\int_{u_{i\infty}^{\mathrm{can}}}\omega
   -\int_0^1H(t,z_\infty(t))\,dt \notag\\
 &=\left(-\int_{w_i}\omega-\int_0^1H_i(t,z_i(t))\,dt\right)
   -\int_{u_{i\infty}^{\mathrm{can}}}\omega \notag\\
 &\quad-\left(\int_0^1H(t,z_\infty(t))\,dt
   -\int_0^1H_i(t,z_i(t))\,dt\right) \notag\\
 &=\mathcal A_{H_i}([z_i,w_i])
   -\int_{u_{i\infty}^{\mathrm{can}}}\omega \notag\\
 &\quad-\left(\int_0^1H(t,z_\infty(t))\,dt
   -\int_0^1H_i(t,z_i(t))\,dt\right).
 \label{eq:21-limit-action-computation}
\end{align}
由于 $z_i$ 一致收敛到 $z_\infty$ 且 $H_i\to H$，
\begin{equation}
 -\left(\int_0^1H(t,z_\infty(t))dt
 -\int_0^1H_i(t,z_i(t))dt\right)\longrightarrow0.
 \label{eq:21-hamiltonian-integral-limit}
\end{equation}
结合~\eqref{eq:21-peak-action-limit}--\eqref{eq:21-hamiltonian-integral-limit}，得到
\[
 \lim_{i\to\infty}\mathcal A_H([z_\infty,w'_i])=\rho(H;a).
\]
故 $\mathcal A_H([z_\infty,w'_i])$ 是柯西序列，从而
\[
 \left|\int_{w'_i}\omega-\int_{w'_j}\omega\right|
 =|\mathcal A_H([z_\infty,w'_i])-\mathcal A_H([z_\infty,w'_j])|
 \longrightarrow0,
\]
即
\[
 \int_{w'_i\#\overline{w'_j}}\omega\longrightarrow0.
\]
因 $\Gamma$ 离散且 $[w'_i\#\overline{w'_j}]_\omega\in\Gamma$，对充分大的
$i,j\in\mathbb Z_+$，必有
\begin{equation}
 \int_{w'_i\#\overline{w'_j}}\omega=0.
 \label{eq:21-rational-area-stabilization}
\end{equation}
特别地，数集 $\{\int_{w'_i}\omega\}_{i\in\mathbb Z_+}$ 有界；必要时取子列，可认为
序列 $\int_{w'_i}\omega$ 最终稳定。由 $\mathcal A_H$ 的定义，对任意充分大的 $N$，
$\mathcal A_H([z_\infty,w'_N])=
\lim_{i\to\infty}\mathcal A_H([z_\infty,w'_i])$。回到
\eqref{eq:21-limit-action-computation} 可见作用量本身稳定，故
\[
 \mathcal A_H([z_\infty,w'_N])
 =\lim_{i\to\infty}\mathcal A_H([z_\infty,w'_i])=\rho(H;a).
\]
所以 $\rho(H;a)$ 确为 $\mathcal A_H$ 在临界点 $[z_\infty,w'_N]$ 处的值。
这就是对无理情形不成立的关键所在。
\end{proof}

上述有理情形证明中，下述引理起关键作用。部分由于它在无理情形下失效，即使只论
非退化情形，无理情形谱性的证明也远比有理情形非平凡。

\begin{lemma}
\label{lem:21-rational-critical-set-compactness}
当 $(M,\omega)$ 有理时，对任意光滑哈密顿量 $K$，
$\operatorname{Crit}\mathcal A_K\subset\widetilde{\mathcal L}_0(M)$ 是
$\mathbb R$ 的闭子集，并且在覆盖空间
$\pi:\widetilde{\mathcal L}_0(M)\to\mathcal L_0(M)$ 的子空间拓扑中局部紧。
\footnote{译者注：原文称
$\operatorname{Crit}\mathcal A_K\subset\widetilde{\mathcal L}_0(M)$ 是
“$\mathbb R$ 的闭子集”，但证明只说明它在覆盖空间中闭；译文照录该类型疑点。}
\end{lemma}

\begin{proof}
当 $(M,\omega)$ 有理时，上述 $\pi$ 的覆盖群 $\Gamma$ 离散。再结合紧流形 $M$ 上
方程 $\dot z=X_K(z)$ 的解集紧，可知
\[
 \operatorname{Crit}(\mathcal A_K)
 =\{[z,w]\in\widetilde{\mathcal L}_0(M)\mid\dot z=X_K(z)\}
\]
是闭集且局部紧。
\end{proof}

无理情形非退化谱性的证明最早由 (Oh09a) 给出（半正情形）。该证明使用连续性方法，
从一个小莫尔斯函数出发沿\emph{瑟夫同伦}推进，并细致分析与谱不变量对应的临界值如何
\stfhIndex{I190}{Eshe}{厄舍（Usher）}
变化。后来厄舍 (Ush08) 基于非阿基米德分析中一条优美的逼近定理，在 13.8 节所述
弗洛尔同调公理框架中给出简洁的代数证明。这里只陈述厄舍的第 13.8 节框架下
非退化谱性公理的抽象表述；其证明
依赖紧闭链的存在性，细节见 (Ush08)。回忆定义 13.8.1 中的
$\operatorname{Spec}(\mathfrak c)$。

\begin{theorem}[定理 1.4 (Ush08)]
\label{thm:21-usher-spectrality}
设 $\mathfrak c$ 是诺特环 $R$ 上的滤过弗洛尔--诺维科夫复形。则对每个
$a\in H_*(\mathfrak c)$，存在 $\alpha\in C_*(\mathfrak c)$ 使
$[\alpha]=a$ 且 $\rho(a)=\lambda(\alpha)$。特别地，
$\rho(a)\in\operatorname{Spec}(\mathfrak c)$。
\end{theorem}

对任意非退化哈密顿量 $H$ 的弗洛尔复形
$(CF_*(H),\partial_{(H,J)})$，此定理说明谱数 $\rho(H;a)$ 由一条紧弗洛尔闭链
$\alpha$ 的层级实现，即 $\lambda_H(\alpha)=\rho(H;a)$。特别地，
$\rho(H;a)\in\operatorname{Spec}(H)$，这就证明了 $\rho(H;a)$ 的非退化谱性。
